\chapter{Quick Reference}\label{ch:reference}

\chapterguide%
  {Everything here has appeared earlier with full explanation; nothing in this appendix is new material.}
  {revise before an assessment, pick a method for a new problem, and check a property claim quickly.}

This appendix gathers the decisions and numbers spread across the previous chapters. Use it
for revision rather than first learning: every entry points back to the chapter that explains
it.

\section{Choosing a search algorithm}

The five search algorithms in Chapters~\ref{ch:uninformed}--\ref{ch:searchpatterns} differ in
one respect only --- what the frontier is sorted on --- so the choice reduces to a short chain
of questions about the problem (\figref{fig:choose-search}).

\begin{aifig}[Choosing a search algorithm]{A decision path for the five search algorithms in
this course. Each question is about the \emph{problem}, not the code: whether the steps cost
the same, and whether a usable estimate of remaining distance exists.}{fig:choose-search}
\begin{tikzpicture}[font=\small,
  q/.style={draw=MLOrange,fill=MLPaleOrange,rounded corners=2mm,text width=3.7cm,
            align=center,inner sep=2.2mm,font=\scriptsize\color{MLNavy}},
  ans/.style={draw=MLGreen,fill=MLPaleTeal,rounded corners=2mm,text width=3.0cm,
            align=center,inner sep=2.2mm,font=\scriptsize\bfseries\color{MLNavy}},
  yes/.style={aiarrow,draw=MLGreen},
  no/.style={aiarrow,draw=MLRed}]

\node[q] (q1) at (0,0) {Do you have a usable estimate $h(n)$ of the distance remaining?};
\node[q] (q2) at (0,-2.5) {Do all steps cost the same?};
\node[q] (q3) at (6.0,0) {Is the estimate \emph{admissible} (never overestimates)?};
\node[ans] (astar) at (11.4,0.6) {A*\\{\scriptsize\mdseries optimal, \chapref{ch:informed}}};
\node[ans] (gbfs)  at (11.4,-1.1) {GBFS\\{\scriptsize\mdseries fast, not optimal}};
\node[ans] (bfs)   at (6.0,-2.5) {BFS\\{\scriptsize\mdseries \chapref{ch:uninformed}}};
\node[ans] (ucs)   at (6.0,-4.2) {UCS\\{\scriptsize\mdseries \chapref{ch:uninformed}}};

\draw[yes] (q1)--(q3) node[midway,above,ailabel,color=MLGreen] {yes};
\draw[no]  (q1)--(q2) node[midway,left,ailabel,color=MLRed] {no};
\draw[yes] (q3)--(astar) node[midway,above,ailabel,color=MLGreen] {yes};
\draw[no]  (q3)--(gbfs) node[midway,below,ailabel,color=MLRed] {no};
\draw[yes] (q2)--(bfs) node[midway,above,ailabel,color=MLGreen] {yes};
\draw[no]  (q2.south) |- (ucs.west) node[pos=0.25,left,ailabel,color=MLRed] {no};

\node[ailabel,anchor=north west,align=left,text width=5.0cm] at (-1.9,-5.3)
  {DFS does not appear on this path because it is never the choice when a
   \emph{good} path is wanted. Use it when memory is the binding constraint,
   or when any reachable solution will do.};
\end{tikzpicture}
\end{aifig}

\section{Search algorithm properties}

\begin{mltable}
\begin{tabularx}{\linewidth}{@{}>{\bfseries\leavevmode\color{MLNavy}}p{0.10\linewidth}
  p{0.13\linewidth}p{0.15\linewidth}p{0.16\linewidth}p{0.12\linewidth}X@{}}
\toprule
\rowcolor{MLPaleBlue!92}
Method & Sorts on & Complete? & Optimal? & Space & Use when \\
\midrule
BFS   & depth      & yes, finite $b$ & only if equal costs & $O(b^{d})$ & fewest steps wanted \\
DFS   & recency    & no, in general  & no                  & $O(bm)$    & memory is tight \\
UCS   & $g(n)$     & yes, costs $\ge\varepsilon$ & yes, non-negative costs & large & weighted, no heuristic \\
GBFS  & $h(n)$     & no              & no                  & large      & speed matters more than cost \\
A*    & $g+h$      & yes             & yes, if $h$ admissible & large   & weighted, heuristic available \\
\bottomrule
\end{tabularx}
\end{mltable}

\begin{keyidea}
Read the space column first. BFS, UCS, GBFS, and A* all store the whole frontier; on a large
problem it is memory, not time, that stops them. DFS is the only one that does not, and it
pays for that with both completeness and optimality.
\end{keyidea}

\section{Fuzzy logic, Mamdani, Sugeno, and ANFIS}

\begin{mltable}
\begin{tabularx}{\linewidth}{@{}>{\bfseries\leavevmode\color{MLNavy}}p{0.15\linewidth}
  p{0.22\linewidth}p{0.22\linewidth}X@{}}
\toprule
\rowcolor{MLPaleBlue!92}
& Mamdani & Sugeno & ANFIS \\
\midrule
Output form  & fuzzy set, then defuzzified & constant or linear function & Sugeno only \\
Rules from   & an expert & an expert & learned from data \\
MF shapes from & an expert & an expert & learned from data \\
Best for     & interpretable control & optimisation, computation & data-driven modelling \\
Needs data?  & no & no & yes, and enough of it \\
\bottomrule
\end{tabularx}
\end{mltable}

\section{The Lab 7 NLP methods by mechanism}

\begin{mltable}
\begin{tabularx}{\linewidth}{@{}>{\bfseries\leavevmode\color{MLNavy}}p{0.24\linewidth}
  p{0.26\linewidth}X@{}}
\toprule
\rowcolor{MLPaleBlue!92}
Mechanism & Examples & Fails when \\
\midrule
Rules and dictionaries & NER, grammar correction, lexicon sentiment, language detection
  & the input contains anything not in the list \\
Text statistics & frequency summarisation, TF--IDF, keyword extraction
  & the corpus is too small for the counts to mean anything \\
Trained model & Naive Bayes sentiment
  & the training data does not resemble the input \\
\bottomrule
\end{tabularx}
\end{mltable}

\section{Genetic algorithm settings}

\begin{mltable}
\begin{tabularx}{\linewidth}{@{}>{\bfseries\leavevmode\color{MLNavy}}p{0.20\linewidth}
  p{0.20\linewidth}X@{}}
\toprule
\rowcolor{MLPaleBlue!92}
Setting & Typical effect of raising it & Symptom that it is wrong \\
\midrule
Population size   & better coverage, slower generations & premature convergence if too small \\
Mutation rate     & more exploration, less stability & fitness oscillates and never settles \\
Crossover rate    & faster mixing of good material & diversity collapses if selection is also strong \\
Generations       & more refinement & best-so-far curve flat for many generations \\
Elitism           & guarantees the best is never lost & population converges on one answer early \\
\bottomrule
\end{tabularx}
\end{mltable}

\section{Mistakes that recur across the labs}

\begin{mltable}
\begin{tabularx}{\linewidth}{@{}>{\bfseries\leavevmode\color{MLNavy}}p{0.30\linewidth}X@{}}
\toprule
\rowcolor{MLPaleBlue!92}
Mistake & What to do instead \\
\midrule
Marking BFS states visited on removal & Mark on insertion, or the same state enters the queue repeatedly \\
Assuming a heuristic is admissible & Compare $h(n)$ against the true $h^\star(n)$ before trusting A* (\chapref{ch:informed}) \\
Expecting BFS to return the cheapest path & BFS counts edges, not cost (\figref{fig:ucs-vs-bfs}) \\
Normalising fitness that can be negative & Shift or transform scores before dividing (\figref{fig:roulette}) \\
Updating \texttt{prev\_error} too early in PID & The derivative term becomes zero (\chapref{ch:robotics}) \\
Judging a controller from one input & Sweep the whole range and check the shape (\chapref{ch:fuzzy}) \\
Scaling train and test data separately & Fit the scaler on training data only, then apply it \\
\bottomrule
\end{tabularx}
\end{mltable}
